\section*{Capítulo \ref{ch:b2:isotope-geochemistry} --- Geoquímica isotópica}

\begin{solution}{exo:b2:isotope-geochemistry:1}
Tristão: $(0.51249/0.512638 - 1)\times10^4 = -2.9$. Mangaia: $(0.51285/0.512638 -
1)\times10^4 = +4.1$. MORB: $+8.5$. As fontes de Mangaia e do MORB, com
$\varepsilon_{\mathrm{Nd}}$ positivo, tiveram em média uma Sm/Nd mais alta que a dos \omterm{def:b2:planet-formation-meteorites:chondrule}{condritos}; a
fonte de Tristão, uma mais baixa.
\end{solution}

\begin{solution}{exo:b2:isotope-geochemistry:2}
$\delta^{18}\mathrm{O}_{\mathrm{VPDB}} = (\num{2.0672e-3}/\num{2.0052e-3} - 1)\times1000 =
+30.9$ por mil na escala \omterm{def:b2:isotope-geochemistry:standards}{VSMOW}. A calcita tem $^{18}\mathrm{O}/^{16}\mathrm{O} =
\num{2.0672e-3}\times(1 - 0.0020) = \num{2.0631e-3}$, logo $\delta^{18}\mathrm{O} =
(\num{2.0631e-3}/\num{2.0052e-3} - 1)\times1000 = +28.9$ por mil (\omterm{def:b2:isotope-geochemistry:standards}{VSMOW}).
\end{solution}

\begin{solution}{exo:b2:isotope-geochemistry:3}
A \qty{1.0}{Ga}: $e^{\lambda t} - 1 = 0.006545$, e $0.512638 - 0.1966\times0.006545 =
0.511351$. A \qty{4.567}{Ga}: $e^{\lambda t}-1 = 0.030243$, e $0.512638 -
0.1966\times0.030243 = 0.506692$.
\end{solution}

\begin{solution}{exo:b2:isotope-geochemistry:4}
$\delta\mathrm{D} = 8\times(-18.0) + 10 = -134$ por mil. Para o lago, $d = -42 -
8\times(-4.0) = -10$ por mil, abaixo dos 10 da reta: a água fica abaixo da \omterm{prop:b2:isotope-geochemistry:mwl}{reta
meteórica}, como as águas superficiais evaporadas (os isótopos leves saíram primeiro).
\end{solution}

\begin{solution}{exo:b2:isotope-geochemistry:5}
$Q = 10^4\times\num{6.524e-12}\times0.1966/0.512638 = \num{2.502e-8}$ por ano, ou
\qty{25.0}{\per\giga\year}. O granito tem $\mathrm{Sm}/\mathrm{Nd} = 6.0/30 = 0.200$,
contra $0.145/0.457 = 0.317$ para o CI: $f = 0.200/0.317 - 1 = -0.37$. Ao longo de
\qty{0.5}{Gyr}: $\Delta\varepsilon = 25.0\times(-0.37)\times0.5 = -4.6$ unidades.
\end{solution}

\begin{solution}{exo:b2:isotope-geochemistry:6}
$\varepsilon \approx 3.5 + 25.0\times(-0.45)\times2.7 = 3.5 - 30.4 = -26.9$, contra
$-27.3$: a aproximação é 0.4 unidade alta demais (não negativa o bastante). A maior parte da
diferença, cerca de 0.3 unidade, vem de $e^{\lambda t} > 1$: o pai, e portanto o
déficit de crescimento em relação ao CHUR, era mais abundante no passado que hoje
(\cref{exo:b2:isotope-geochemistry:10}); o resto vem do arredondamento de $f$ e do valor inicial.
\end{solution}

\begin{solution}{exo:b2:isotope-geochemistry:7}
Com uma fração em massa $x$ de granito, $\mathrm{Sr} = 320x + 130(1-x)$ e
$R = (320x\times0.7200 + 130(1-x)\times0.7030)/\mathrm{Sr}$. Para $x = 0.05$: Sr
\qty{139.5}{ppm}, $R = 0.70495$; $x = 0.10$: 149.0, 0.70665; $x = 0.20$: 168.0,
0.70948. Pela demonstração da proposição de mistura, $R\,\mathrm{Sr}$ é linear em Sr:
$R = 0.731632 - 3.7221/\mathrm{Sr}$, uma reta de inclinação $-3.722$ (ppm) e
intercepto 0.7316 em função de $1/\mathrm{Sr}$. Verificação em $x = 0.05$: $0.731632 - 3.7221/139.5
= 0.70495$.
\end{solution}

\begin{solution}{exo:b2:isotope-geochemistry:8}
A \qty{20}{\degreeCelsius}: $\delta^{18}\mathrm{O} = 1000/1.0098 - 1000 = -9.70$ por
mil, $\delta\mathrm{D} = 1000/1.084 - 1000 = -77.5$ por mil, inclinação $77.5/9.70 = 7.99$.
A \qty{0}{\degreeCelsius}: $-11.56$ e $-99.9$ por mil, inclinação 8.64. Ambas estão próximas
da inclinação 8 da \omterm{prop:b2:isotope-geochemistry:mwl}{reta meteórica}: só o fracionamento de equilíbrio da água
entre líquido e vapor já fixa perto de 8 a razão entre os deslocamentos do deutério e do
oxigênio-18.
\end{solution}

\begin{solution}{exo:b2:isotope-geochemistry:9}
A crosta contém $0.0051\times27/1.25 = 0.110$, 11 por cento do neodímio. Balanço
de massa: $0.110\times(-15) + 0.890\,\varepsilon_m = 0$, logo $\varepsilon_m = 1.65/0.890 =
+1.9$. Com um reservatório empobrecido contendo uma fração $x$ a $+8.5$ e o resto a 0:
$0.110\times(-15) + 8.5x = 0$, $x = 1.65/8.5 = 0.19$: cerca de um quinto do
neodímio.
\end{solution}

\begin{solution}{exo:b2:isotope-geochemistry:10}
A taxa de crescimento é $\lambda p\,e^{\lambda t}$, com $e^{\lambda t} \approx 1 + \lambda t$;
em média sobre o intervalo que vai da separação $T$ até o presente, o fator é
$1 + \lambda T/2$, de modo que os termos desprezados têm tamanho relativo $\lambda T/2$. Ao longo
da idade da Terra, $\num{6.524e-12}\times\num{4.567e9}/2 = 0.015$, 1.5 por cento. Para a
crosta do \cref{exo:b2:isotope-geochemistry:6}, $\lambda T/2 = 0.0088$, e o erro é $0.0088\times30.4 = 0.27$
unidade. A lenta variação da razão do CHUR pela qual se divide $\varepsilon$ acrescenta
outro termo, menor, da ordem de $(p/R)\lambda T \times \varepsilon_0$, aqui 0.02 unidade.
\end{solution}

\begin{solution}{exo:b2:isotope-geochemistry:11}
$K = 13.3/5.0 = 2.66$. Com $s = 0.5$: $s' = 2.66\times0.5/(1 + 1.66\times0.5) =
1.33/1.83 = 0.73$. A diferença $s' - s$ é máxima onde
$\mathrm{d}s'/\mathrm{d}s = K/(1+(K-1)s)^2 = 1$, em $s = (\sqrt K - 1)/(K-1) =
0.631/1.66 = 0.38$, onde $s' = 0.62$: uma diferença de 0.24.
\end{solution}

\begin{solution}{exo:b2:isotope-geochemistry:12}
Geração: $2.3$ a $2.6\times\qty{7.2e9}{km^3}/\qty{4.5e9}{yr} = 3.7$ a
\qty{4.2}{km^3} por ano. Destruição: $1.3$ a $1.6\times7.2/4.5 = 2.1$ a
\qty{2.6}{km^3} por ano. As taxas atuais diferem: os acréscimos (2.7) estão abaixo da
geração média do modelo, a remoção (3.2) acima de sua destruição média. No modelo, a crosta cresceu
cedo, quando a geração era mais rápida, e o presente é uma época de equilíbrio ou de leve
perda, como diz o balanço das zonas de subducção.
\end{solution}

\begin{solution}{pb:b2:isotope-geochemistry:1}
\textbf{1.} Manto: $(0.51307/0.512638 - 1)\times10^4 = +8.4$. Sedimento:
$(0.51141/0.512638 - 1)\times10^4 = -24.0$.

\textbf{2.} Kora $+2.8$, Vel $-4.3$, Asha $+4.7$.

\textbf{3.} Com $^{87}\mathrm{Sr}/^{86}\mathrm{Sr}$ na horizontal e $\varepsilon_{\mathrm{Nd}}$ na vertical:
manto (0.70282, $+8.4$), Kora (0.70410, $+2.8$), Vel (0.70618, $-4.3$), sedimento
(0.72000, $-24.0$), Asha (0.70290, $+4.7$). Só Vel, com um
$\varepsilon_{\mathrm{Nd}}$ negativo, fica do lado enriquecido do CHUR.

\textbf{4.} Abaixo da faixa: na razão de estrôncio do \omterm{def:b2:isotope-geochemistry:reservoirs}{manto empobrecido}, a faixa
teria $\varepsilon_{\mathrm{Nd}}$ perto de $+8$; Asha tem $+4.7$. Ela se parece com o
grupo de Santa Helena, de tipo HIMU (Mangaia, Tubuai, Santa Helena).

\textbf{5.} A fusão altera as razões pai--filho, mas não, de modo mensurável, as
razões isotópicas do estrôncio e do neodímio (proposição sobre a fusão); as lavas são
jovens, de modo que suas razões não cresceram desde a erupção (questão 6).

\textbf{6.} $\Delta R = 0.12\,(e^{\lambda\cdot10^6} - 1) = 0.12\times\num{6.52e-6} =
\num{7.8e-7}$, ou $\num{7.8e-7}/0.512638\times10^4 = 0.015$ unidade épsilon: muito abaixo das
diferenças entre as ilhas, e abaixo da precisão de uma medida.

\textbf{7.} $\mathrm{Nd} = 40f + 0.75(1-f)$, e
$R = \bigl(40f\times0.51141 + 0.75(1-f)\times0.51307\bigr)/\mathrm{Nd}$.

\textbf{8.} $f = 0.5$ por cento: Nd \qty{0.9463}{ppm}, $\varepsilon = +1.6$; 1 por cento:
1.1425, $-2.9$; 2 por cento: 1.535, $-8.4$; 5 por cento: 2.7125, $-15.4$.

\textbf{9.} Sr 10.95, 11.90, 13.80 e \qty{19.50}{ppm}; $^{87}\mathrm{Sr}/^{86}\mathrm{Sr} =
0.70439$, 0.70571, 0.70780 e 0.71163.

\textbf{10.} $R\,\mathrm{Sr} = 10\times0.70282 + f(200\times0.72000 - 10\times0.70282)$ e
$\mathrm{Sr} = 10 + 190f$; eliminando $f$, $R\,\mathrm{Sr} = 7.0282 +
(\mathrm{Sr}-10)\times136.9718/190$, logo $R = 0.720904 - 0.18084/\mathrm{Sr}$: uma
reta de inclinação $-0.1808$ (ppm) em função de $1/\mathrm{Sr}$. Verificação em $f = 1$ por cento:
$0.720904 - 0.18084/11.90 = 0.70571$.

\textbf{11.} $K = (10/0.75)/(200/40) = 13.33/5.0 = 2.67$. O manto é relativamente
mais rico em estrôncio, o sedimento em neodímio, de modo que a razão do Nd se move em direção ao
sedimento mais depressa: um pouco de sedimento desloca muito o $\varepsilon_{\mathrm{Nd}}$ antes que
a $^{87}\mathrm{Sr}/^{86}\mathrm{Sr}$ se mova muito.

\textbf{12.} 53 vezes mais rico em Nd ($40/0.75$), 20 vezes em Sr. Com $f = 1$ por cento, o
sedimento fornece $0.40/1.1425 = 35$ por cento do neodímio, o que desloca
$\varepsilon_{\mathrm{Nd}}$ em $0.35\times(8.4 + 24.0) = 11.3$ unidades, mas só $2.0/11.9 =
17$ por cento do estrôncio.

\textbf{13.} $0.01\times40/1.1425 = 0.350$: 35 por cento.

\textbf{14.} De $R\,(40f + 0.75(1-f)) = 40fR_S + 0.75(1-f)R_D$:
$f = 0.75(R - R_D)/\bigl(40(R_S - R) + 0.75(R - R_D)\bigr)$, e o mesmo com 200
e 10 para o estrôncio.

\textbf{15.} $f = 0.75\times(-0.00029)/\bigl(40\times(-0.00137) + 0.75\times(-0.00029)\bigr)
= \num{-2.175e-4}/\num{-5.502e-2} = 0.0040$: 0.40 por cento.

\textbf{16.} $f = 10\times0.00128/\bigl(200\times0.01590 + 10\times0.00128\bigr) =
0.0128/3.1928 = 0.0040$: 0.40 por cento; os dois concordam.

\textbf{17.} Neodímio: $0.75\times(-0.00065)/\bigl(40\times(-0.00101) +
0.75\times(-0.00065)\bigr) = \num{-4.875e-4}/\num{-4.089e-2} = 0.0119$; estrôncio:
$10\times0.00336/\bigl(200\times0.01382 + 10\times0.00336\bigr) = 0.0336/2.7976 = 0.0120$.
Ambos dão 1.2 por cento.

\textbf{18.} $f = 0.75\times(-0.00019)/\bigl(40\times(-0.00147) + 0.75\times(-0.00019)\bigr)
= 0.0024$, o que prevê Sr \qty{10.46}{ppm} e $^{87}\mathrm{Sr}/^{86}\mathrm{Sr} =
0.70361$, contra 0.70290 medido. Asha não está sobre a curva de mistura manto--sedimento:
sua fonte contém outro componente.

\textbf{19.} Crosta oceânica reciclada antiga, o tipo HIMU: seus isótopos de Nd exigem uma Sm/Nd
um pouco abaixo da do \omterm{def:b2:isotope-geochemistry:reservoirs}{manto empobrecido}, seus isótopos de Sr uma Rb/Sr baixa, e seu
chumbo (não medido aqui) uma U/Pb alta, tudo isso durante um a dois bilhões de anos.

\textbf{20.} $e^{\lambda_{147}\cdot2\times10^9} - 1 = 0.013133$. Sedimento: $0.51141 -
0.110\times0.013133 = 0.509965$; CHUR: $0.512638 - 0.1966\times0.013133 = 0.510056$;
$\varepsilon_{\mathrm{Nd}}(2\,\mathrm{Ga}) = (0.509965/0.510056 - 1)\times10^4 = -1.8$.

\textbf{21.} $f = 0.110/0.1966 - 1 = -0.440$; desde a subducção, $\varepsilon$ variou
$25.0\times(-0.440)\times2.0 = -22.0$, logo $\varepsilon(2\,\mathrm{Ga}) \approx -24.0 + 22.0
= -2.0$, a 0.2 unidade do valor exato.

\textbf{22.} $e^{\lambda_{87}\cdot2\times10^9} - 1 = 0.028338$;
$0.72000 - 0.50\times0.028338 = 0.70583$.

\textbf{23.} Na subducção, o sedimento tinha $\varepsilon_{\mathrm{Nd}} = -1.8$, próximo dos
$-4.7$ de uma crosta com apenas 0.7 bilhão de anos naquela época, e bem acima dos
valores fortemente negativos da crosta antiga: ele corresponde à erosão de crosta jovem, talvez
misturada com algum material derivado do manto. Sua $^{87}\mathrm{Sr}/^{86}\mathrm{Sr}$ de
0.7058 naquela época também é a de uma rocha crustal jovem.

\textbf{24.} A fonte de Vel, a ilha mais enriquecida, contém
\textbf{\qty[text-series-to-math=true]{1.2}{\percent}} de sedimento reciclado.
\end{solution}
